\begin{abstract}
\noindent
Internet of Things (IoT) devices generate network traffic whose interpretation depends on both individual connections and the communication patterns around them. Conventional machine learning can classify connections from numerical measurements, while deep learning can learn richer representations of those measurements. Graph neural networks (GNNs) add communication context by representing devices as nodes and connections as edges. Pretrained language models provide another view by encoding textual descriptions of network events. Earlier work has combined graph and language models through parallel fusion, explanations and iterative feedback. These approaches motivate testing whether their interaction improves edge-level intrusion detection. This report presents GLASS (Graph-Language Adaptive Semantic-Structural fusion), an edge-level intrusion-detection pipeline combining a graph attention network with a frozen cybersecurity language encoder. GLASS describes each aggregated flow using a fixed text template, combines graph and language embeddings through a gate that varies across edges and feature dimensions, and evaluates an iterative feedback stage that consults a semantic classifier on uncertain edges. The evaluation compares a prototype semantic scorer and a trained classification head, including the head alone, using pooled out-of-fold predictions across three training seeds. Experimental results on the NF-UNSW-NB15 and NF-ToN-IoT NetFlow datasets show that the trained-head feedback configuration improves over the GNN, prototype semantic and fusion stages by statistically separated margins on both datasets. However, the trained head alone has slightly higher mean macro-F1, and its difference from feedback is not separated. Fusion improves over the prototype semantic model on both datasets, but does not establish an improvement over the GNN. With output fusion disabled at the earlier consultation settings tested, neither realistic consultant provides a separated injection gain, and trained-head injection produces a separated loss on NF-ToN-IoT; an oracle supplied with true labels improves through the same channel. GLASS therefore provides evidence for the value of a trained semantic classifier and a controlled way to examine graph--language interaction, while an additional benefit from feedback remains unestablished. The findings concern label-conditioned aggregated graphs and do not estimate deployment performance.
\end{abstract}

\vspace{1em}
\noindent\textbf{Keywords:} GLASS, intrusion detection, graph neural networks, edge classification,
structural--semantic fusion, uncertainty-guided feedback, negative results.
